\documentclass[aspectratio=169,10pt]{beamer}

% Shared Beamer setup for Fall 2026 course slides.
% Clean Cal Poly-inspired theme: no custom class files or uncommon packages.
\mode<presentation>{
  \usetheme{default}
  \usefonttheme{professionalfonts}
}

\definecolor{CalPolyGreen}{HTML}{154734}
\definecolor{CalPolyGold}{HTML}{BD8B13}
\definecolor{CalPolySage}{HTML}{7FA28F}
\definecolor{CalPolyMint}{HTML}{DDF1E7}
\definecolor{CalPolySand}{HTML}{A29F86}
\definecolor{CalPolyBeige}{HTML}{EFEEDE}
\definecolor{DeckInk}{HTML}{1F2933}
\definecolor{DeckMuted}{HTML}{5F6B7A}
\definecolor{DeckSoft}{HTML}{F7F7F5}

\setbeamersize{text margin left=0.55in,text margin right=0.55in}
\setbeamercolor{normal text}{fg=DeckInk,bg=white}
\setbeamercolor{structure}{fg=CalPolyGreen}
\setbeamercolor{title}{fg=white}
\setbeamercolor{frametitle}{fg=CalPolyGreen,bg=white}
\setbeamercolor{block title}{bg=CalPolySage,fg=white}
\setbeamercolor{block body}{bg=CalPolyMint,fg=black}
\setbeamercolor{alerted text}{fg=CalPolyGold}

\setbeamerfont{title}{size=\Huge,series=\bfseries}
\setbeamerfont{frametitle}{size=\Large,series=\bfseries}
\setbeamerfont{block title}{size=\normalsize,series=\bfseries}

\setbeamertemplate{navigation symbols}{}
\setbeamertemplate{itemize item}{\color{CalPolyGold}$\blacktriangleright$}
\setbeamertemplate{itemize subitem}{\color{CalPolyGreen}$\bullet$}
\setbeamertemplate{footline}{
  \hfill{\color{DeckMuted}\scriptsize\insertframenumber/\inserttotalframenumber}
  \hspace{0.18in}\vspace{0.12in}
}
\setbeamertemplate{frametitle}{
  \vspace{0.18in}
  \hspace{0.02in}\insertframetitle\par
  \vspace{0.08in}
  \noindent\makebox[\linewidth][l]{%
    {\color{CalPolyGold}\rule{0.55in}{2pt}}%
    {\color{CalPolyGreen}\rule{\dimexpr\linewidth-0.55in\relax}{2pt}}%
  }
}

\newcommand{\CourseNumber}{Course}
\newcommand{\CourseTitle}{Course Title}
\newcommand{\LectureNumber}{0}
\newcommand{\LectureTitle}{Lecture Title}
\newcommand{\LectureDate}{Date}
\newcommand{\CourseUnit}{Unit}

\newcommand{\coursenumber}[1]{\renewcommand{\CourseNumber}{#1}}
\newcommand{\coursetitle}[1]{\renewcommand{\CourseTitle}{#1}}
\newcommand{\lecturenumber}[1]{\renewcommand{\LectureNumber}{#1}}
\newcommand{\lecturetitle}[1]{\renewcommand{\LectureTitle}{#1}}
\newcommand{\lecturedate}[1]{\renewcommand{\LectureDate}{#1}}
\newcommand{\courseunit}[1]{\renewcommand{\CourseUnit}{#1}}

\author{Kirk Duran}
\institute{California Polytechnic State University, San Luis Obispo}
\date{\LectureDate}

\newcommand{\makedecktitle}{
  \begingroup
  \setbeamercolor{background canvas}{bg=CalPolyGreen}
  \begin{frame}[plain]
    \vfill
    \begin{center}
      {\usebeamerfont{title}\color{white}\LectureTitle\par}
      \vspace{0.18in}
      {\Large\color{white}Lecture \LectureNumber\par}
    \end{center}
    \vfill
    \noindent{\color{white}\rule{\linewidth}{1.2pt}}\par
    \vspace{0.08in}
    \noindent\colorbox{CalPolyGold}{\makebox[0.24\linewidth][c]{\color{white}\Large\CourseNumber}}
    \hspace{0.12in}{\color{white}\Large Kirk Duran}
  \end{frame}
  \endgroup
}

\newenvironment{keyidea}{\begin{block}{Key idea}}{\end{block}}
\newenvironment{activity}{\begin{block}{In class}}{\end{block}}
\newenvironment{cpdefinition}{\begin{block}{Definition}}{\end{block}}
\newenvironment{exampleblockcp}{
  \begingroup
  \setbeamercolor{block title}{bg=CalPolySand,fg=white}
  \setbeamercolor{block body}{bg=CalPolyBeige,fg=black}
  \begin{block}{Example}
}{
  \end{block}
  \endgroup
}

\newcommand{\imageplaceholder}[2][0.3in]{%
  \noindent\makebox[\linewidth][c]{%
    \fcolorbox{CalPolySage}{CalPolyBeige}{%
      \parbox[c][#1][c]{0.86\linewidth}{%
        \centering
        {\color{DeckMuted}\scriptsize Image placeholder: }%
        {\color{DeckInk}\scriptsize\ttfamily\detokenize{#2}}%
      }%
    }%
  }%
}

\newcommand{\sectionframe}[1]{
  \begingroup
  \setbeamercolor{background canvas}{bg=CalPolyGreen}
  \begin{frame}[plain]
    \vfill
    {\color{white}\LARGE\bfseries #1\par}
    \vspace{0.18in}
    {\color{CalPolyGold}\rule{0.22\paperwidth}{3pt}}
    {\color{white}\rule{0.62\paperwidth}{3pt}}
    \vfill
  \end{frame}
  \endgroup
}

\newcommand{\placeholdernote}{
  \begin{block}{Placeholder}
    Replace these bullets with the final lecture material, examples, and in-class activities.
  \end{block}
}


\usepackage{tikz}

\graphicspath{{../assets/lecture-15/}}

% If an image file is not yet available, skip the visual slot.
\newcommand{\lectureimage}[2][1.75in]{%
  \IfFileExists{../assets/lecture-15/#2}{%
    \includegraphics[width=\linewidth,height=#1,keepaspectratio]{#2}%
  }{}%
}

% Explicit placeholder for visuals/examples that will be added later.
\newcommand{\lectureplaceholder}[2][2.35in]{%
  \begin{center}
    \fbox{%
      \begin{minipage}[c][#1][c]{0.90\linewidth}
        \centering
        \small\textit{#2}
      \end{minipage}%
    }
  \end{center}%
}

\setbeamersize{text margin left=0.42in,text margin right=0.42in}

\coursenumber{CS 4880}
\coursetitle{Artificial Intelligence}
\lecturenumber{15}
\lecturetitle{Automated Logical Inference}
\lecturedate{September 30, 2026}
\courseunit{Logic and Knowledge Representation}

\begin{document}

% Estimated pacing: 50 minutes
% Motivation + bridge from model checking: 6 min
% CNF: 11 min
% Resolution: 15 min
% Horn clauses + forward chaining: 9 min
% Backward chaining + comparison: 7 min
% Wrap / bridge to FOL: 2 min

% -----------------------------------------------------------------------------
% Slide 1
% -----------------------------------------------------------------------------
\makedecktitle

\section{From Semantics to Algorithms}

% -----------------------------------------------------------------------------
% Slide 2
% -----------------------------------------------------------------------------
\begin{frame}[shrink=10]{Last Time: We Could Prove Things by Checking Every World}
  \begin{columns}[T,onlytextwidth]
    \begin{column}{0.55\textwidth}
      Suppose our knowledge base contains
      \[
        A \rightarrow B,\qquad A.
      \]

      To test whether
      \[
        KB \models B,
      \]
      truth-table inference checks every assignment of truth values.

      \begin{block}{Correct, but expensive}
        With $n$ propositional symbols, there are
        \[
          2^n
        \]
        possible models.
      \end{block}
    \end{column}
    \begin{column}{0.39\textwidth}
      \lectureimage[2.70in]{slide2.png}
    \end{column}
  \end{columns}

  \begin{keyidea}
    Semantic inference asks whether a conclusion is true in every model. Today we ask whether we can derive conclusions without enumerating all models.
  \end{keyidea}
\end{frame}

% -----------------------------------------------------------------------------
% Slide 3
% -----------------------------------------------------------------------------
\begin{frame}[shrink=10]{A Logical Agent Needs an Inference Engine}
  \begin{columns}[T,onlytextwidth]
    \begin{column}{0.53\textwidth}
      A logical agent stores sentences in a knowledge base, but storing knowledge is not enough.

      The agent also needs procedures that can answer questions such as:
      \begin{itemize}
        \item Is this conclusion implied by what I know?
        \item Which new facts can I derive?
        \item Can I prove a particular goal?
      \end{itemize}

      \begin{block}{Today}
        We turn logical inference into algorithms that manipulate sentences directly.
      \end{block}
    \end{column}
    \begin{column}{0.41\textwidth}
      \lectureimage[2.75in]{slide3.png}
    \end{column}
  \end{columns}
\end{frame}

% -----------------------------------------------------------------------------
% Slide 4
% -----------------------------------------------------------------------------
\begin{frame}[shrink=10]{Two Ways to Reason}
  \begin{columns}[T,onlytextwidth]
    \begin{column}{0.47\textwidth}
      \begin{block}{Model checking}
        Ask whether the knowledge base and query have the right truth values across possible worlds.

        \vspace{0.08in}
        \textbf{Think:} inspect models.
      \end{block}
    \end{column}
    \begin{column}{0.47\textwidth}
      \begin{block}{Theorem proving}
        Apply inference rules that transform known sentences into new sentences.

        \vspace{0.08in}
        \textbf{Think:} manipulate symbols.
      \end{block}
    \end{column}
  \end{columns}

  \lectureimage[1.45in]{slide4.png}

  \begin{keyidea}
    A proof system is useful when a small sequence of syntactic steps can replace an enormous semantic search.
  \end{keyidea}
\end{frame}

\section{Conjunctive Normal Form}

% -----------------------------------------------------------------------------
% Slide 5
% -----------------------------------------------------------------------------
\begin{frame}[shrink=10]{Why Normalize Logical Sentences?}
  \begin{columns}[T,onlytextwidth]
    \begin{column}{0.54\textwidth}
      Consider two equivalent statements:
      \[
        A \rightarrow B
      \]
      and
      \[
        \neg A \lor B.
      \]

      A human can recognize that these mean the same thing. An inference algorithm is easier to write if equivalent sentences are converted into a predictable structure first.

      \begin{block}{Normalization}
        Rewrite many different logical expressions into one standard form before inference begins.
      \end{block}
    \end{column}
    \begin{column}{0.40\textwidth}
      \lectureimage[2.55in]{slide5.png}
    \end{column}
  \end{columns}
\end{frame}

% -----------------------------------------------------------------------------
% Slide 6
% -----------------------------------------------------------------------------
\begin{frame}[shrink=10]{Literals and Clauses}
  \begin{cpdefinition}
    A \textbf{literal} is either a propositional symbol or its negation.
  \end{cpdefinition}

  \begin{columns}[T,onlytextwidth]
    \begin{column}{0.46\textwidth}
      \begin{block}{Literals}
        \[
          A,\qquad \neg B,\qquad C
        \]
      \end{block}
    \end{column}
    \begin{column}{0.48\textwidth}
      \begin{block}{Clause}
        A disjunction of literals:
        \[
          A \lor \neg B \lor C
        \]
      \end{block}
    \end{column}
  \end{columns}

  \begin{keyidea}
    Resolution operates on clauses, so clause structure is the basic unit of automated inference in this lecture.
  \end{keyidea}
\end{frame}

% -----------------------------------------------------------------------------
% Slide 7
% -----------------------------------------------------------------------------
\begin{frame}[shrink=10]{Conjunctive Normal Form}
  \begin{cpdefinition}
    A sentence is in \textbf{conjunctive normal form (CNF)} when it is a conjunction of clauses, where each clause is a disjunction of literals.
  \end{cpdefinition}

  \[
    (A \lor \neg B) \land (C \lor D) \land (\neg A \lor E)
  \]

  \begin{center}
    \textbf{AND of ORs}
  \end{center}

  \lectureimage[1.55in]{slide7.png}
\end{frame}

% -----------------------------------------------------------------------------
% Slide 8
% -----------------------------------------------------------------------------
\begin{frame}[shrink=10]{CNF Conversion: The Recipe}
  A common propositional conversion pipeline is:

  \begin{enumerate}
    \item Eliminate biconditionals $\leftrightarrow$.
    \item Eliminate implications $\rightarrow$.
    \item Move negations inward using De Morgan's laws and double negation.
    \item Distribute $\lor$ over $\land$.
  \end{enumerate}

  \begin{block}{Important}
    Each rewrite preserves logical equivalence. We are changing the sentence's form, not its meaning.
  \end{block}

  \lectureimage[1.30in]{slide8.png}
\end{frame}

% -----------------------------------------------------------------------------
% Slide 9
% -----------------------------------------------------------------------------
\begin{frame}[shrink=8]{Worked CNF Conversion: Start with an Implication}
  Convert
  \[
    A \rightarrow (B \land C)
  \]
  to CNF.

  \begin{block}{Step 1: eliminate implication}
    \[
      \neg A \lor (B \land C)
    \]
  \end{block}

  We are not done because the OR contains a conjunction.

  \begin{center}
    \textbf{Next: distribute $\lor$ over $\land$.}
  \end{center}
\end{frame}

% -----------------------------------------------------------------------------
% Slide 10
% -----------------------------------------------------------------------------
\begin{frame}[shrink=8]{Worked CNF Conversion: Distribute}
  Continue from
  \[
    \neg A \lor (B \land C).
  \]

  Distribution gives
  \[
    (\neg A \lor B) \land (\neg A \lor C).
  \]

  Now the expression is a conjunction of clauses.

  \begin{block}{Clause set view}
    \[
      \{\neg A,B\},\qquad \{\neg A,C\}
    \]
  \end{block}

  \begin{keyidea}
    CNF lets the inference engine treat the knowledge base as a collection of clauses instead of arbitrary nested expressions.
  \end{keyidea}
\end{frame}

% -----------------------------------------------------------------------------
% Slide 11
% -----------------------------------------------------------------------------
\begin{frame}[shrink=10]{Quick Check: Which One Is CNF?}
  Which of the following is already in conjunctive normal form?

  \begin{enumerate}
    \item $(A \land B) \lor C$
    \item $(A \lor B) \land (\neg C \lor D)$
    \item $A \rightarrow (B \lor C)$
  \end{enumerate}

  \pause
  \begin{block}{Answer}
    \[
      (A \lor B) \land (\neg C \lor D)
    \]
    is an AND of OR-clauses.
  \end{block}
\end{frame}

\section{Resolution}

% -----------------------------------------------------------------------------
% Slide 12
% -----------------------------------------------------------------------------
\begin{frame}[shrink=10]{Resolution: One Small Rule with a Big Job}
  Suppose the knowledge base contains
  \[
    A \lor B
  \]
  and
  \[
    \neg B \lor C.
  \]

  If $B$ is true, the second clause requires $C$. If $B$ is false, the first clause requires $A$.

  Either way,
  \[
    A \lor C
  \]
  must be true.

  \begin{keyidea}
    Resolution removes a pair of complementary literals and combines what remains.
  \end{keyidea}
\end{frame}

% -----------------------------------------------------------------------------
% Slide 13
% -----------------------------------------------------------------------------
\begin{frame}[shrink=8]{The Resolution Rule}
  \begin{cpdefinition}
    If two clauses contain complementary literals, resolution derives a new clause containing the remaining literals.
  \end{cpdefinition}

  \[
    (L \lor A_1 \lor \cdots \lor A_m)
  \]
  \[
    (\neg L \lor B_1 \lor \cdots \lor B_n)
  \]

  resolve to
  \[
    A_1 \lor \cdots \lor A_m \lor B_1 \lor \cdots \lor B_n.
  \]

  \lectureimage[1.35in]{slide13.png}
\end{frame}

% -----------------------------------------------------------------------------
% Slide 14
% -----------------------------------------------------------------------------
\begin{frame}[shrink=8]{Worked Resolution Step}
  Start with
  \[
    \neg Rain \lor Wet
  \]
  and
  \[
    Rain.
  \]

  Treat the second statement as the one-literal clause $(Rain)$.

  Resolve on the complementary literals $Rain$ and $\neg Rain$.

  Removing that complementary pair leaves
  \[
    Wet.
  \]

  \begin{block}{This should look familiar}
    Resolution has reproduced modus ponens using a general clause-level rule.
  \end{block}
\end{frame}

% -----------------------------------------------------------------------------
% Slide 15
% -----------------------------------------------------------------------------
\begin{frame}[shrink=10]{But How Do We Prove a Query?}
  Suppose we want to determine whether
  \[
    KB \models Q.
  \]

  Resolution is usually used by \textbf{refutation}:

  \begin{enumerate}
    \item Add the negation of the query, $\neg Q$, to the knowledge base.
    \item Repeatedly apply resolution.
    \item If this produces a contradiction, then $Q$ must follow from the original knowledge base.
  \end{enumerate}

  \begin{center}
    \textbf{Assume the conclusion is false and see whether the knowledge base becomes impossible.}
  \end{center}
\end{frame}

% -----------------------------------------------------------------------------
% Slide 16
% -----------------------------------------------------------------------------
\begin{frame}[shrink=10]{Resolution by Refutation: Running Example}
  Knowledge base:
  \[
    A \rightarrow B
  \]
  \[
    B \rightarrow C
  \]
  \[
    A
  \]

  Query:
  \[
    C?
  \]

  Convert the knowledge base to clauses:
  \[
    \neg A \lor B
  \]
  \[
    \neg B \lor C
  \]
  \[
    A
  \]

  Then add the negated query:
  \[
    \neg C.
  \]
\end{frame}

% -----------------------------------------------------------------------------
% Slide 17
% -----------------------------------------------------------------------------
\begin{frame}[shrink=8]{Refutation Setup: Clause Form}
  The proof search begins with four clauses:

  \[
    \neg A \lor B
  \]
  \[
    \neg B \lor C
  \]
  \[
    A
  \]
  \[
    \neg C \qquad \text{(assumed for contradiction)}
  \]

  \begin{keyidea}
    From this point forward, every inference step uses the same resolution rule.
  \end{keyidea}
\end{frame}

% -----------------------------------------------------------------------------
% Slide 18
% -----------------------------------------------------------------------------
\begin{frame}[shrink=8]{Refutation Trace: Resolve the Query Backward}
  Start with
  \[
    \neg B \lor C
  \]
  and the assumed negated query
  \[
    \neg C.
  \]

  Resolving on $C$ gives
  \[
    \neg B.
  \]

  \begin{block}{Interpretation}
    If $C$ were false, then the rule $B\rightarrow C$ says $B$ would also have to be false.
  \end{block}

  \lectureimage[1.25in]{slide18.png}
\end{frame}

% -----------------------------------------------------------------------------
% Slide 18
% -----------------------------------------------------------------------------
\begin{frame}[shrink=8]{Refutation Trace: Continue the Chain}
  Now resolve
  \[
    \neg A \lor B
  \]
  with
  \[
    \neg B.
  \]

  Resolving on $B$ gives
  \[
    \neg A.
  \]

  But the knowledge base also contains
  \[
    A.
  \]

  The assumption $\neg C$ is forcing us toward a direct contradiction.
\end{frame}

% -----------------------------------------------------------------------------
% Slide 19
% -----------------------------------------------------------------------------
\begin{frame}[shrink=8]{The Empty Clause Means Contradiction}
  Resolve
  \[
    A
  \]
  with
  \[
    \neg A.
  \]

  Nothing remains:
  \[
    \Box
  \]

  The \textbf{empty clause} represents an impossible clause: a disjunction with no literal that could make it true.

  \begin{block}{Conclusion}
    $KB \land \neg C$ is unsatisfiable, therefore
    \[
      KB \models C.
    \]
  \end{block}
\end{frame}

% -----------------------------------------------------------------------------
% Slide 20
% -----------------------------------------------------------------------------
\begin{frame}[shrink=10]{Why Refutation Works}
  To prove
  \[
    KB \models Q,
  \]
  we ask whether there exists any model in which
  \[
    KB \land \neg Q
  \]
  is true.

  If no such model exists, then every model satisfying $KB$ must satisfy $Q$.

  \begin{keyidea}
    Resolution is syntactic, but refutation connects it directly back to the semantic definition of entailment from the previous lecture.
  \end{keyidea}
\end{frame}

% -----------------------------------------------------------------------------
% Slide 21
% -----------------------------------------------------------------------------
\begin{frame}[shrink=10]{Resolution Is Sound and Complete for Propositional Logic}
  \begin{columns}[T,onlytextwidth]
    \begin{column}{0.47\textwidth}
      \begin{block}{Sound}
        Every clause derived by resolution is logically implied by its parent clauses.

        It does not invent false consequences.
      \end{block}
    \end{column}
    \begin{column}{0.47\textwidth}
      \begin{block}{Complete for refutation}
        If a propositional CNF knowledge base is unsatisfiable, resolution can derive the empty clause.

        Given enough search, a contradiction can be found.
      \end{block}
    \end{column}
  \end{columns}

  \begin{alertblock}{But complete does not mean cheap}
    The number of possible clauses and resolution steps can still grow explosively.
  \end{alertblock}
\end{frame}

% -----------------------------------------------------------------------------
% Slide 22
% -----------------------------------------------------------------------------
\begin{frame}[shrink=10]{Resolution Is Also a Search Problem}
  At each step, an inference engine must decide:
  \begin{itemize}
    \item Which two clauses should be resolved?
    \item Which complementary literal should be used?
    \item Which derived clauses should be kept?
    \item When should the search stop?
  \end{itemize}

  \lectureimage[1.55in]{slide23.png}

  \begin{keyidea}
    Logical inference does not escape search. It changes the search space from world states to proofs and symbolic consequences.
  \end{keyidea}
\end{frame}

% -----------------------------------------------------------------------------
% Slide 23
% -----------------------------------------------------------------------------
\begin{frame}[shrink=10]{A Useful Restriction: Horn Clauses}
  General CNF is expressive, but unrestricted resolution can be expensive.

  A \textbf{Horn clause} contains at most one positive literal.

  Examples:
  \[
    \neg A \lor \neg B \lor C
  \]
  \[
    \neg A \lor D
  \]
  \[
    E
  \]

  These have a convenient rule interpretation.
\end{frame}

\section{Horn Clauses and Rule-Based Inference}

% -----------------------------------------------------------------------------
% Slide 24
% -----------------------------------------------------------------------------
\begin{frame}[shrink=8]{Horn Clauses Look Like Rules}
  Consider
  \[
    \neg A \lor \neg B \lor C.
  \]

  This is equivalent to
  \[
    (A \land B) \rightarrow C.
  \]

  We can read it as
  \begin{center}
    \textbf{IF $A$ and $B$, THEN $C$.}
  \end{center}

  \begin{block}{Rule-system vocabulary}
    $A$ and $B$ are \textbf{premises}.\qquad $C$ is the \textbf{conclusion} or \textbf{head}.
  \end{block}

  \lectureimage[1.20in]{slide25.png}
\end{frame}

% -----------------------------------------------------------------------------
% Slide 25
% -----------------------------------------------------------------------------
\begin{frame}[shrink=10]{Why Horn Clauses Matter in AI}
  Horn rules support a simple and efficient style of inference.

  \begin{itemize}
    \item Facts can be stored directly.
    \item Rules have a natural IF--THEN interpretation.
    \item New facts can be generated by matching satisfied premises.
    \item The same rule representation supports both forward and backward reasoning.
  \end{itemize}

  \begin{block}{Historical connection}
    Rule-based expert systems and logic-programming languages made this style of symbolic reasoning a central part of classical AI.
  \end{block}
\end{frame}

% -----------------------------------------------------------------------------
% Slide 26
% -----------------------------------------------------------------------------
\begin{frame}[shrink=10]{Forward Chaining: Start with What You Know}
  \begin{cpdefinition}
    \textbf{Forward chaining} repeatedly applies rules whose premises are known to be true, adding their conclusions to the knowledge base.
  \end{cpdefinition}

  \begin{center}
    \textbf{Facts $\rightarrow$ applicable rules $\rightarrow$ new facts $\rightarrow$ more rules}
  \end{center}

  \lectureimage[1.65in]{slide27.png}

  \begin{keyidea}
    Forward chaining is data-driven: begin with available evidence and ask what follows from it.
  \end{keyidea}
\end{frame}

% -----------------------------------------------------------------------------
% Slide 27
% -----------------------------------------------------------------------------
\begin{frame}[shrink=8]{Forward Chaining Example: A Maintenance Agent}
  Initial facts:
  \[
    BatteryLow,\qquad ChargerConnected,\qquad Docked
  \]

  Rules:
  \[
    BatteryLow \land ChargerConnected \rightarrow Charging
  \]
  \[
    Charging \land Docked \rightarrow SafeToPowerDown
  \]
  \[
    SafeToPowerDown \rightarrow MaintenanceMode
  \]

  \begin{block}{Question}
    What new facts will forward chaining derive?
  \end{block}
\end{frame}

% -----------------------------------------------------------------------------
% Slide 28
% -----------------------------------------------------------------------------
\begin{frame}[shrink=8]{Forward Chaining Trace}
  \begin{enumerate}
    \item $BatteryLow$ and $ChargerConnected$ are known.
    \[
      \Rightarrow Charging
    \]

    \item $Charging$ and $Docked$ are now known.
    \[
      \Rightarrow SafeToPowerDown
    \]

    \item $SafeToPowerDown$ is known.
    \[
      \Rightarrow MaintenanceMode
    \]
  \end{enumerate}

  \begin{block}{Result}
    The agent derives $MaintenanceMode$ even though that fact was never explicitly entered into the original knowledge base.
  \end{block}
\end{frame}

% -----------------------------------------------------------------------------
% Slide 29
% -----------------------------------------------------------------------------
\begin{frame}[shrink=10]{An Efficient Forward-Chaining Implementation}
  A practical implementation avoids rescanning every rule from scratch.

  \begin{itemize}
    \item Maintain an \textbf{agenda} of newly discovered facts.
    \item For each rule, track how many premises are still unsatisfied.
    \item When the count reaches zero, add the rule's conclusion to the agenda.
    \item Never process the same inferred fact twice.
  \end{itemize}

  \begin{keyidea}
    For propositional Horn knowledge bases, forward chaining can be implemented in time linear in the size of the knowledge base.
  \end{keyidea}
\end{frame}

% -----------------------------------------------------------------------------
% Slide 30
% -----------------------------------------------------------------------------
\begin{frame}[shrink=10]{Backward Chaining: Start with What You Want to Prove}
  \begin{cpdefinition}
    \textbf{Backward chaining} begins with a query and recursively asks which rules could establish that query.
  \end{cpdefinition}

  \begin{center}
    \textbf{Goal $\rightarrow$ rules that conclude the goal $\rightarrow$ required subgoals}
  \end{center}

  \lectureimage[1.65in]{slide31.png}

  \begin{keyidea}
    Backward chaining is goal-driven: explore only the parts of the knowledge base relevant to the current query.
  \end{keyidea}
\end{frame}

% -----------------------------------------------------------------------------
% Slide 31
% -----------------------------------------------------------------------------
\begin{frame}[shrink=8]{Backward Chaining Example}
  Query:
  \[
    MaintenanceMode?
  \]

  The only rule that can produce it is
  \[
    SafeToPowerDown \rightarrow MaintenanceMode.
  \]

  So backward chaining creates the subgoal
  \[
    SafeToPowerDown?
  \]

  To prove that, it finds
  \[
    Charging \land Docked \rightarrow SafeToPowerDown,
  \]
  producing two more subgoals:
  \[
    Charging?\qquad Docked?
  \]
\end{frame}

% -----------------------------------------------------------------------------
% Slide 32
% -----------------------------------------------------------------------------
\begin{frame}[shrink=8]{Backward Chaining Continues Until It Hits Facts}
  We already know
  \[
    Docked.
  \]

  To prove $Charging$, use
  \[
    BatteryLow \land ChargerConnected \rightarrow Charging.
  \]

  Both premises are known facts:
  \[
    BatteryLow,\qquad ChargerConnected.
  \]

  Therefore the proof succeeds all the way back up:
  \[
    MaintenanceMode.
  \]

  \lectureimage[1.20in]{slide33.png}
\end{frame}

% -----------------------------------------------------------------------------
% Slide 33
% -----------------------------------------------------------------------------
\begin{frame}[shrink=10]{Forward vs. Backward Chaining}
  \begin{center}
    \renewcommand{\arraystretch}{1.3}
    \begin{tabular}{l|p{0.34\textwidth}|p{0.34\textwidth}}
      & \textbf{Forward chaining} & \textbf{Backward chaining} \\\hline
      Starts from & known facts & a query / goal \\
      Direction & data $\rightarrow$ consequences & goal $\rightarrow$ prerequisites \\
      Good when & many consequences may matter & one specific query matters \\
      Risk & derive many irrelevant facts & revisit subgoals / chase unproductive rules \\
    \end{tabular}
  \end{center}

  \begin{keyidea}
    The knowledge representation may be the same; the search strategy through the rule space is different.
  \end{keyidea}
\end{frame}

% -----------------------------------------------------------------------------
% Slide 34
% -----------------------------------------------------------------------------
\begin{frame}[shrink=10]{Logical Agents Can Mix Inference Strategies}
  A practical logical agent does not need one universal inference procedure.

  \begin{itemize}
    \item Use forward chaining to maintain obvious consequences of new sensor data.
    \item Use backward chaining when answering a specific query.
    \item Use resolution when general CNF reasoning is required.
    \item Use specialized algorithms when the knowledge representation permits them.
  \end{itemize}

  \begin{block}{AI engineering lesson}
    Representation and inference are tightly coupled. A restricted representation can make reasoning dramatically cheaper.
  \end{block}
\end{frame}

% -----------------------------------------------------------------------------
% Slide 35
% -----------------------------------------------------------------------------
\begin{frame}[shrink=10]{The GOFAI Connection: Intelligence as Proof and Rules}
  Much of classical symbolic AI was built around the idea that intelligent behavior could emerge from:

  \begin{center}
    \textbf{explicit knowledge + symbolic rules + automated inference}
  \end{center}

  This produced theorem provers, expert systems, planning systems, and logic-programming environments.

  \begin{columns}[T,onlytextwidth]
    \begin{column}{0.53\textwidth}
      \begin{block}{Strength}
        Reasoning steps can be explicit, inspectable, and justified.
      \end{block}
    \end{column}
    \begin{column}{0.41\textwidth}
      \begin{block}{Limitation}
        Real-world knowledge can be difficult to encode completely and brittle when assumptions change.
      \end{block}
    \end{column}
  \end{columns}

  \lectureimage[1.00in]{slide36.png}
\end{frame}

% -----------------------------------------------------------------------------
% Slide 36
% -----------------------------------------------------------------------------
\begin{frame}[shrink=8]{Putting the Lecture Together}
  \begin{center}
    \begin{tabular}{c}
      Arbitrary propositional sentences \\
      $\Downarrow$ \\
      \textbf{CNF: standard clause representation} \\
      $\Downarrow$ \\
      \textbf{Resolution: general clause inference} \\
      $\Downarrow$ \\
      \textbf{Horn restriction: rule-like structure} \\
      $\Downarrow$ \\
      \textbf{Forward / backward chaining: efficient rule inference}
    \end{tabular}
  \end{center}

  \begin{keyidea}
    Automated inference is the search for consequences. The structure of the knowledge base determines which inference algorithms are available and how expensive that search becomes.
  \end{keyidea}
\end{frame}

% -----------------------------------------------------------------------------
% Slide 37
% -----------------------------------------------------------------------------
\begin{frame}[shrink=8]{Where Propositional Logic Hits a Wall}
  Suppose we want to represent:

  \begin{center}
    \textbf{Every student enrolled in AI has access to the AI lab.}
  \end{center}

  In propositional logic, we would need separate statements for every student:
  \[
    EnrolledAlice \rightarrow AccessAlice
  \]
  \[
    EnrolledBob \rightarrow AccessBob
  \]
  \[
    EnrolledCarol \rightarrow AccessCarol
  \]
  \[
    \vdots
  \]

  \begin{block}{Next lecture}
    First-order logic introduces objects, predicates, relations, and quantifiers so one rule can describe an entire category of entities.
  \end{block}
\end{frame}

% -----------------------------------------------------------------------------
% Slide 38
% -----------------------------------------------------------------------------
\begin{frame}[shrink=8]{Takeaways and Next Time}
  \begin{itemize}
    \item CNF rewrites propositional sentences into a conjunction of clauses.
    \item Resolution combines clauses by eliminating complementary literals.
    \item Resolution refutation proves $KB\models Q$ by deriving a contradiction from $KB\land\neg Q$.
    \item Horn clauses provide a restricted rule-like form that supports efficient inference.
    \item Forward chaining is data-driven; backward chaining is goal-driven.
    \item Logical inference is still a search problem---but now the search is over proofs, rules, and consequences.
    \item Next: \textbf{First-Order Logic and Knowledge Representation}.
  \end{itemize}

  \begin{keyidea}
    The central tradeoff is expressive power versus inference cost: richer representations can say more, but often require more sophisticated reasoning.
  \end{keyidea}
\end{frame}

\end{document}
